% !TeX root = ../../Thesis.tex
\begin{figure}[htbp]
  \centering
  \begin{tikzpicture}[
    every node/.style={font=\small},
    unit/.style={
      draw, rounded corners=2pt, line width=.7pt,
      minimum width=20mm, minimum height=8mm
    },
    master/.style={unit, draw=hpcgreen, fill=hpcgreen!8},
    worker/.style={unit, draw=hpcblue, fill=hpcblue!8},
    wire/.style={draw=hpcgray, line width=.7pt},
    feed/.style={wire, ->, >=stealth}
  ]
    \node[master] (sched) {Scheduler};
    \node[worker, below=11mm of sched] (n2) {Node 2};
    \node[worker, left=7mm of n2] (n1) {Node 1};
    \node[worker, right=7mm of n2] (n3) {Node 3};

    % Vertical trunk, horizontal bus, then one drop per node.
    \coordinate (bus) at ($(sched.south)!.5!(n2.north)$);
    \draw[wire] (sched.south) -- (bus);
    \draw[wire] (n1.north |- bus) -- (n3.north |- bus);
    \foreach \n in {n1,n2,n3}{
      \draw[feed] (\n.north |- bus) -- (\n.north);
    }
  \end{tikzpicture}
  \caption{\ThesisText{Example architecture diagram drawn with TikZ: a scheduler distributing work to three compute nodes.}{Beispielhaftes Architekturdiagramm in TikZ: ein Scheduler verteilt Arbeit auf drei Rechenknoten.}}
  \label{fig:example-architecture}
\end{figure}
